\section{Prefix function, KMP}
\whenuse{Buscar exactamente un patr\'on en un texto o analizar bordes y per\'iodos
de una cadena sin hashing.}

\begin{lstlisting}[style=cpp]
vector<int> prefix_function(const string &s) {
    int n = s.size();
    vector<int> pi(n, 0);
    for (int i = 1; i < n; ++i) {
        int j = pi[i - 1];
        while (j > 0 && s[i] != s[j]) j = pi[j - 1];
        if (s[i] == s[j]) ++j;
        pi[i] = j;
    }
    return pi;
}

vector<int> kmp_occurrences(const string &text, const string &pattern) {
    if (pattern.empty()) {
        vector<int> all(text.size() + 1);
        iota(all.begin(), all.end(), 0);
        return all;
    }
    vector<int> pi = prefix_function(pattern), answer;
    int matched = 0;
    for (int i = 0; i < (int)text.size(); ++i) {
        while (matched > 0 && text[i] != pattern[matched])
            matched = pi[matched - 1];
        if (text[i] == pattern[matched]) ++matched;
        if (matched == (int)pattern.size()) {
            answer.push_back(i - matched + 1);
            matched = pi[matched - 1];
        }
    }
    return answer;
}
\end{lstlisting}

\complexity{$O(n+m)$ tiempo; $O(m)$ memoria para buscar.}
\subsection{Ejemplo de KMP}

\texttt{pi[i]} es la longitud del mayor prefijo propio de
\texttt{s[0..i]} que también es sufijo. KMP usa esta información para no volver
al inicio del patrón después de una diferencia.

\begin{lstlisting}[style=cpp]
string text="bananana",pattern="ana";
vector<int> positions=kmp_occurrences(text,pattern);
for(int position:positions) cout<<position<<' '; // 1 3 5
\end{lstlisting}

También permite encontrar el mayor borde de una cadena:

\begin{lstlisting}[style=cpp]
string s="abcab";
vector<int> pi=prefix_function(s);
int longest_border=pi.back(); // 2: "ab"
\end{lstlisting}

Un borde es un prefijo propio que también aparece como sufijo.
\section{Z-function}
\whenuse{Conocer para cada posici\'on cu\'anto coincide el sufijo que empieza all\'i
con el prefijo; \`util para patrones, bordes y periodicidad.}

\texttt{z[i]} es la longitud del prefijo de $s$ que coincide desde $i$.

\begin{lstlisting}[style=cpp]
vector<int> z_function(const string &s) {
    int n = s.size();
    vector<int> z(n, 0);
    for (int i = 1, l = 0, r = 0; i < n; ++i) {
        if (i < r) z[i] = min(r - i, z[i - l]);
        while (i + z[i] < n && s[z[i]] == s[i + z[i]]) ++z[i];
        if (i + z[i] > r) l = i, r = i + z[i];
    }
    return z;
}
\end{lstlisting}

\complexity{$O(n)$ tiempo y memoria.}

\subsection{Ejemplo de Z-function}

\texttt{z[i]} indica cuántos caracteres desde \texttt{s[i]} coinciden con el
prefijo completo de \texttt{s}. Por ejemplo, para \texttt{``aaaaa''}:

\[
z=[0,4,3,2,1].
\]

Para buscar un patrón, concatenar patrón, un separador ausente y texto:

\begin{lstlisting}[style=cpp]
string text="bananana",pattern="ana";
string combined=pattern+"#"+text;
vector<int> z=z_function(combined);
int offset=pattern.size()+1;
for(int i=offset;i<(int)combined.size();++i)
    if(z[i]>=(int)pattern.size())
        cout<<i-offset<<' '; // 1 3 5
\end{lstlisting}

\warning{El separador, aquí \texttt{\#}, no debe aparecer ni en el patrón ni en
el texto.}

\begin{itemize}
    \item \textbf{KMP/prefix function:} buscar exactamente un patrón y estudiar
    bordes o períodos.
    \item \textbf{Z-function:} conocer cuánto coincide cada sufijo con el prefijo;
    suele resultar más directa para matching mediante concatenación.
    \item Ambos son deterministas y cuestan $O(n+m)$; no tienen colisiones como
    rolling hash.
\end{itemize}